\clearpage
\VerdierPhysicalPageAnchor{0236}%
\begingroup\small\setlength{\parskip}{0pt}\setlength{\emergencystretch}{1.5em}
\setlength{\abovedisplayskip}{4pt plus 2pt minus 2pt}%
\setlength{\belowdisplayskip}{4pt plus 2pt minus 2pt}%

\VerdierSemanticAnchor{ch4d-4-2-4}
\noindent\textbf{4.2.4.}\quad
$\operatorname{comp}(\mathcal A)$ の $\widetilde{\mathbb Z}$ 型フィルター付き
対象の射は、命題 (4.2.3) a) で述べた対応するスペクトル対象の間の射を
明らかな仕方で誘導する。したがって命題 (4.2.3) は、
$\operatorname{comp}(\mathcal A)$ のフィルター付き対象の圏から、
$D(\mathcal A)$ に値をとるスペクトル対象の圏への関手を定める。

\VerdierSemanticAnchor{ch4d-4-2-4-first-canonical}
$X$ を $\operatorname{comp}(\mathcal A)$ の対象とし、
$X(q)\subset X$、$q\in\mathbb Z$ をその第一標準フィルトレーション
とする。この場合、命題 (4.2.3) がフィルター付き対象
$\bigl(X,X(q),q\in\mathbb Z\bigr)$ に対応させるスペクトル対象は、複体
$X$ に付随する、$D(\mathcal A)$ に値をとる第一スペクトル対象 (4.1.5)
にほかならない。

\VerdierSemanticAnchor{ch4d-4-2-4-nfold}
同様に、$X$ を $\mathcal A$ の $n$-重複複体、$\int_n X=Y$ を付随する
単純複体とし、$m$ を整数、$1\leq m\leq n$ とする。また
\[
0={}^{m}\!Y(-\infty)\subset\cdots\subset{}^{m}\!Y(q)
\subset\cdots\subset{}^{m}\!Y(+\infty)=Y
\]
を、$X$ の第 $m$ 次数により定まるフィルトレーション (4.1.6.9) とする。
この場合、命題 (4.2.3) がフィルター付き対象
$\bigl(Y,{}^{m}\!Y(q),q\in\mathbb Z\bigr)$ に対応させるスペクトル対象は、
第 $m$ 次数でフィルター付けされた $X$ によって定まる、$D(\mathcal A)$
に値をとるスペクトル対象 (4.1.8) にほかならない。

\VerdierSemanticAnchor{ch4d-4-3}
\medskip\noindent\textbf{4.3. 第二標準スペクトル対象。}

\VerdierSemanticAnchor{ch4d-4-3-1}
\noindent\textbf{4.3.1.}\quad
$X$ を $\operatorname{comp}(\mathcal A)$ の対象とする。次のようにおく：

\VerdierSemanticAnchor{ch4d-4-3-1-1}
\begin{equation}
\begin{aligned}
X(-\infty)&=0,\qquad X(+\infty)=X,\\
X(q)&=\cdots\longrightarrow X^{q-3}\longrightarrow X^{q-2}
\longrightarrow\operatorname{Ker}(d_X^{q-1})
\longrightarrow0\longrightarrow0\longrightarrow\cdots,
\qquad q\in\mathbb Z,
\end{aligned}
\tag{4.3.1.1}
\end{equation}
後者の複体の微分は $X$ の微分から誘導される。標準モノ射

\VerdierSemanticAnchor{ch4d-4-3-1-2}
\begin{equation}
0=X(-\infty)\subset\cdots\subset X(q)\subset X(q+1)
\subset\cdots\subset X(+\infty)=X
\tag{4.3.1.2}
\end{equation}
があり、これらは $X$ 上に $\widetilde{\mathbb Z}$ 型フィルトレーションを
定める。これを $X$ の\emph{第二標準フィルトレーション}という。命題
(4.2.3) によりスペクトル対象 $\bigl(X(p,q),\delta(p,q,r)\bigr)$ を得る。これを複体
$X$ に付随する、$D(\mathcal A)$ に値をとる\emph{第二スペクトル対象}という。

\VerdierSemanticAnchor{ch4d-4-3-1-associated-functor}
$\operatorname{comp}(\mathcal A)$ の射 $X\longrightarrow Y$ は第二標準
フィルトレーションを保つので、関手
\[
X\longmapsto\bigl(X(p,q),\delta(p,q,r)\bigr)
\]
は $\operatorname{comp}(\mathcal A)$ から $D(\mathcal A)$ に値をとる
スペクトル対象の圏への関手である。この関手を\emph{第二標準スペクトル対象関手}
という。
\endgroup
